% GeoMERIT - BDCC academic-language revision.
% Built on the official MDPI LaTeX template (Definitions/mdpi.cls).
% Language and argumentation revised against the supplied reference-paper corpus.
% Section 3 figures updated: Figure 1 now combines the new well-location map,
% the long-tail class distribution panel, and the new family-grouped missingness matrix.
\documentclass[bdcc,article,submit,pdftex,moreauthors]{Definitions/mdpi}

% Compatibility for MiKTeX installations where colortbl is newer than array.
\makeatletter
\providecommand\insert@pcolumn{\insert@column}
\makeatother

\firstpage{1}
\makeatletter
\setcounter{page}{\@firstpage}
\makeatother
\pubvolume{1}
\issuenum{1}
\articlenumber{0}
\pubyear{2026}
\copyrightyear{2026}
\datereceived{ }
\daterevised{ }
\dateaccepted{ }
\datepublished{ }

\Title{GeoMERIT: Missingness-Aware and Geological-Cost-Guided Lithology Prediction across Field Wells}

\newcommand{\orcidauthorA}{0000-0000-0000-000X}
\Author{Anonymous Author $^{1,}$*\orcidA{}}
\AuthorNames{Anonymous Author}
\address{$^{1}$ \quad Affiliation 1; email@example.com}
\corres{Correspondence: email@example.com}

\abstract{Robust lithology prediction from well logs is complicated by structured variation
in curve availability and acquisition conditions, severe class imbalance, label ambiguity
near formation boundaries, and asymmetric geological misclassification costs. These factors
can produce optimistic estimates under random row-wise splits and reduce predictive stability
when a model is applied to unseen wells. We propose GeoMERIT, a missingness-aware and
penalty-guided method that addresses these factors at the representation, learning,
posterior-modeling, and decision stages. The representation encodes missing curves, anomalous
responses, curve-suite identity, within-well normalized responses, and local depth context
without first reconstructing absent measurements. The learning objective combines class,
boundary, and label-confidence weights. Posterior estimates are obtained from a
three-backend gradient-boosted ensemble and refined using a coarse-to-fine geological
hierarchy and rare-class experts. Final labels are selected by minimizing expected geological
cost under the FORCE penalty matrix after long-tail logit adjustment. On the FORCE 2020 dataset
(118 wells, 12 lithology classes, 1{,}431{,}383 valid labeled samples), under strict 10-fold
well-grouped validation that simulates prediction on unseen wells, GeoMERIT reaches
Weighted F1 $=0.7526$, Boundary F1 $=0.4567$, and Penalty $=-0.5886$, improving on the
published RFE+XGBoost reference by $+3.5\%$, $+11.4\%$, and $+6.4\%$ simultaneously, with a
tail-class mean F1 of $0.5061$. The three gains correspond to three distinct facets of
robustness---overall prediction quality, boundary stability, and the geological severity of
errors. The results indicate that jointly modeling data quality, class imbalance, boundary
uncertainty, and geological cost can improve lithology prediction under heterogeneous
well-log conditions.}

\keyword{robust lithology prediction; well logs; missing data; long-tailed learning;
boundary samples; cost-sensitive decision; gradient-boosted decision trees; FORCE 2020}


% Center short figure captions under the corresponding figure without moving detailed explanations into captions.
\newcommand{\ShortCenteredFigCaption}[2]{%
  \refstepcounter{figure}%
  \label{#2}%
  \vspace{0.08cm}%
  \begin{center}
  \small\textbf{Figure~\thefigure.}~#1
  \end{center}%
}

\begin{document}
\sloppy

\section{Introduction}
\label{sec:intro}
Lithology prediction is a fundamental task in reservoir characterization, geological
modeling, and hydrocarbon development evaluation~\cite{rider2011,ellis2007,doveton2014}. Each
well log records formation responses to a suite of physical measurements---gamma ray
(GR), bulk density (RHOB), neutron porosity (NPHI), photoelectric factor (PEF), sonic transit
time (DTC), resistivity (RDEP), and others---along depth, and machine learning methods are
now routinely used to map these curves to lithology labels at
scale~\cite{hall2016facies,merembayev2021,dev2019formation,carvalho2026benchmark}. Many
learning workflows nevertheless rely on standard supervised-classification assumptions, including
complete inputs, reasonably balanced classes, identically distributed training and test samples,
and equally costly errors. These assumptions are often violated in operational well-log
collections.

Curve availability is commonly structured by tool selection, vendor, acquisition period,
and borehole conditions, while calibration and environmental effects introduce inter-well
shifts in measured responses. The associated labels are strongly imbalanced, and samples
near formation boundaries are often ambiguous. In addition, the geological consequence of
an error depends on the true and predicted lithologies. The FORCE 2020
benchmark~\cite{bormann2020force} exhibits all of these conditions in a public
dataset. Among its 1{,}431{,}383 valid labeled samples from 118 Norwegian North Sea wells,
Shale accounts for 61.3\% of labels while Basement contributes 141 rows (0.01\%). This
head-to-tail ratio of roughly $6220{:}1$ coexists with entire log families missing in specific wells
(Figure~\ref{fig:missingness_matrix}). The benchmark also provides a
$12{\times}12$ penalty matrix in which predicting Sandstone/Shale for true Sandstone incurs
a smaller cost than predicting Halite for the same sample. Together, these properties can
reduce cross-well generalization. Random row-wise splits may produce optimistic results because
well-specific acquisition signatures and local stratigraphic continuity leak between training
and test rows~\cite{rosenblatt2024,roberts2017cv,mila2022nndm}. Models evaluated in this way
may subsequently degrade on unseen wells, incomplete curve suites, rare lithologies, and
boundary intervals.

Most existing approaches focus on individual aspects of this problem, leaving their joint
effect less thoroughly studied. Imputation-centered workflows~\cite{vanbuuren2011mice,stekhoven2012missforest,liu2025wellgan,xie2024cgan}
reconstruct plausible curve values but erase the missingness pattern itself, which in well
logs is informative because the decision to run a curve reflects vintage, vendor, borehole
condition, and operational choices that correlate with geology~\cite{chen2026rfe,little2019stats,josse2024missing}.
Class-imbalance remedies re-weight or over-sample rare
classes~\cite{chawla2002smote,lin2017focal,cui2019classbalanced,he2024imbalanced} but generally do not
distinguish boundary samples from formation-interior samples. Deep
architectures can learn rich representations but typically require complete curve suites and
large labeled corpora, and have not demonstrated consistent superiority over tree ensembles
when evaluation is strictly grouped by
well~\cite{wang2024dnnlitho,grinsztajn2022trees,shwartzziv2022tabular,mcelfresh2024benchmark}.
Most published systems also decode by posterior argmax, which does not account for
asymmetric costs even when an official penalty matrix is
available~\cite{elkan2001cost,berger1985decision}. A recent robust-feature-engineering (RFE)
study with quality-aware XGBoost~\cite{chen2026rfe} established the value of treating
missingness as signal at the feature level, raising Weighted F1 from $0.6640$ to $0.7272$ on
FORCE 2020 under well-grouped validation. Building on that result, this study evaluates
whether robustness can be improved further by treating data quality, long-tail structure,
boundary uncertainty, and geological cost jointly in the representation, training objective,
posterior model, and decision rule.

We present GeoMERIT (Missingness-encoding, Ensembling, Robust-features,
Imbalance-aware, Tail-correction), a robust lithology prediction method for incomplete,
heterogeneous, and long-tailed well-log data. GeoMERIT comprises four components that act at
successive stages of the prediction workflow
(Section~\ref{sec:method}, Figure~\ref{fig:framework}, Table~\ref{tab:layers}). Curve absence,
anomalous responses, curve-suite combinations, within-well normalized responses, and local
depth context form a missingness-aware state representation. Joint class, boundary, and
confidence weighting increases the influence of rare lithologies and formation transitions
during training. A three-backend tree ensemble, a coarse-to-fine geological hierarchy, and
dedicated tail experts are then used to estimate class posteriors under severe imbalance.
The final prediction minimizes expected geological cost under the FORCE penalty matrix after
logit adjustment. Well-grouped validation evaluates stability on wells excluded from training
and is used as the experimental protocol rather than as part of the method definition. Under
this protocol, GeoMERIT improves on the RFE+XGBoost reference
simultaneously on Weighted F1 ($0.7272\!\to\!0.7526$, $+3.5\%$), Boundary F1
($0.4100\!\to\!0.4567$, $+11.4\%$), and FORCE Penalty ($-0.6289\!\to\!-0.5886$, $+6.4\%$),
while raising the tail-class mean F1 to $0.5061$.

The main contribution is an integrated prediction framework that treats
missingness, acquisition heterogeneity, long-tailed supervision, boundary ambiguity, and
asymmetric geological cost as explicit modeling considerations. The framework contains four
design components with corresponding empirical evaluations.
\begin{enumerate}
\item \textbf{A missingness-aware robust geologic state representation.} Curve absence,
anomalous responses, curve-suite identity, well-specific normalization, local depth context,
and stratigraphic metadata are encoded in a unified state vector. This representation retains
information about both observed responses and curve availability, and its behavior is evaluated
on feature-deficient wells
(Sections~\ref{sec:repr} and \ref{sec:rq3}).
\item \textbf{A long-tail- and boundary-enhanced learning mechanism.} A multiplicative sample
weight combining effective-number class weighting, boundary up-weighting, and label-confidence
discounting increases the contribution of rare lithologies and formation-transition samples
to the objective. Tail-class and boundary performance are evaluated in
Sections~\ref{sec:learning} and \ref{sec:rq4}.
\item \textbf{Structured posterior lithology modeling.} GeoMERIT geometrically fuses the
posteriors of LightGBM, XGBoost, and CatBoost, combines them with a
coarse-to-fine geological hierarchy over six super-classes, and gates six one-vs-rest tail
experts. These components provide separate estimates for geological groups and tail classes
before posterior combination (Sections~\ref{sec:posterior} and \ref{sec:rq2}).
\item \textbf{A geological-cost-guided decision layer.} Logit adjustment addresses the
long-tailed training prior, and Bayes-risk decoding over the official FORCE penalty matrix
selects the class with the lowest expected geological cost. The decision-layer settings and
their sensitivity are examined in Sections~\ref{sec:decision} and \ref{sec:rq5}.
\end{enumerate}

The remainder of the paper is organized as follows. Section~\ref{sec:related} positions
GeoMERIT against four families of prior work. Section~\ref{sec:data} describes the dataset,
the metrics, and the well-grouped robustness protocol. Section~\ref{sec:method} formalizes
the four robustness layers. Section~\ref{sec:setup} details the experimental setup and
comparison fairness. Section~\ref{sec:results} reports results, moving from overall
held-out-well performance through the sources of the gains to behavior on feature-deficient
wells, rare lithologies, and the decision layer. Section~\ref{sec:discussion} analyzes failure
modes and limitations, and
Section~\ref{sec:conclusion} concludes.

\section{Related Work}
\label{sec:related}
Prior work is reviewed in four groups according to the aspect of well-log prediction it
addresses: classification and benchmarking, missingness and data quality, long-tailed and
boundary-aware learning, and cost-sensitive decision-making.

\subsection{Well-Log Lithology Classification and Benchmarking}
Machine learning lithology classification has progressed from early facies
challenges~\cite{hall2016facies} through systematic algorithm
comparisons~\cite{merembayev2021,dev2019formation} to deep architectures including CNNs,
recurrent models, transformers, diffusion-based augmentation, and multi-scale hierarchical
networks~\cite{wang2024dnnlitho,chen2025geological,zhao2025diffusion,zhan2025mchgeo}.
Interpretability-oriented hybrids such as SHAP-based meta-learning
transformers~\cite{sun2024shapmeta,lundberg2017shap} and refined-lithology
pipelines~\cite{shi2023refined} demonstrate strong within-dataset accuracy. Benchmark studies
make the comparability problem explicit. Published scores are frequently incomparable because
splits, label sets, and metrics differ, and reported performance can swing by 10--20 absolute
F1 points with the validation protocol alone, because the unit of geological generalization
is the well, not the row~\cite{carvalho2026benchmark,rosenblatt2024}. These findings motivate
methods that explicitly account for the data conditions responsible for cross-well
variability. GeoMERIT adopts well-grouped evaluation and incorporates corresponding
mechanisms in its representation, learning, posterior, and decision stages.

\subsection{Missingness and Data-Quality Modeling in Well Logs}
The dominant response to incomplete logs is reconstruction, most commonly through chained-equation
and forest imputation~\cite{vanbuuren2011mice,stekhoven2012missforest} or through generative
imputation with sequence GANs and conditional GANs~\cite{liu2025wellgan,xie2024cgan}. Reconstruction is
appropriate when missingness is ignorable. Well-log missingness, however, is structurally
non-random because whole curves are absent as a result of tool-string choices that correlate
with vintage, operator, and target formation~\cite{chen2026rfe}. Statistical learning theory shows
that supervised learners can remain consistent when missingness is encoded as input rather
than repaired~\cite{little2019stats,josse2024missing}, and the RFE
study~\cite{chen2026rfe} demonstrated the practical value of this position on FORCE 2020.
The present study examines whether explicit missingness modeling can be combined with class
weighting, structured posterior estimation, and cost-sensitive decoding in a unified
framework.

\subsection{Long-Tailed and Boundary-Aware Learning}
General-purpose remedies for class imbalance include synthetic
oversampling~\cite{chawla2002smote}, focal loss~\cite{lin2017focal}, margin-based
losses~\cite{cao2019ldam}, effective-number class weighting~\cite{cui2019classbalanced}, and
prior-corrected logit adjustment~\cite{menon2021logit}; surveys document their
trade-offs~\cite{he2024imbalanced}. Two domain constraints restrict their transfer to well
logs. First, synthetic interpolation of log rows can violate depth continuity and
stratigraphic consistency, generating physically implausible samples; we therefore avoid
synthetic oversampling entirely. Second, lithology boundaries concentrate both label
ambiguity and practical interest, so imbalance handling must not sacrifice boundary
fidelity---a requirement most imbalance methods never examine. GeoMERIT combines
effective-number weighting~\cite{cui2019classbalanced} with explicit boundary up-weighting and
treats logit adjustment~\cite{menon2021logit} as a prior-shift correction applied at decision
time, where its temperature can be tuned against the geological cost directly.

\subsection{Cost-Sensitive and Geology-Aware Decision Making}
Cost-sensitive learning theory establishes that when misclassification costs are asymmetric,
the optimal decision minimizes expected cost rather than maximizing posterior
probability~\cite{elkan2001cost,berger1985decision,andrade2019bayes,zadrozny2002multiclass}.
FORCE 2020 operationalizes geological cost in its penalty matrix, yet most published
systems---including the RFE reference~\cite{chen2026rfe}---still decode by argmax and report
the penalty only as an evaluation metric. Sequence decoders such as Viterbi smoothing and
(semi-)Markov CRFs~\cite{viterbi1967,lafferty2001crf,sarawagi2004semimarkov} optimize label
continuity but not geological cost. To our knowledge, GeoMERIT is the first FORCE 2020 system
whose headline predictions are produced by Bayes-risk decoding over the official penalty
matrix, coupled with logit adjustment in a single, jointly tuned decision layer.

\section{Dataset and Evaluation Protocol}
\label{sec:data}

\subsection{Dataset}
\label{sec:dataset}
We use the public FORCE 2020 lithofacies dataset~\cite{bormann2020force}, which contains
$1{,}431{,}383$ valid labeled samples from 118 wells and 12 lithology classes. Each sample is
indexed by measured depth and may contain up to 20 well-log curves, including GR, SGR, RHOB,
DRHO, NPHI, PEF, DTC, DTS, RDEP, RSHA, RMED, RXO, RMIC, SP, CALI, BS, DCAL, ROP, ROPA, and
MUDWEIGHT. We also use well-head coordinates and NPD stratigraphic tables, including GROUP,
FORMATION, and casing-depth information. These variables define a realistic lithology
prediction setting in which the model must infer rock type from heterogeneous curve suites,
uneven sample support, and geological labels whose errors have unequal consequences.

The first source of heterogeneity is spatial and operational. The 118 wells are distributed
across multiple North Sea blocks, and the available curve suites differ across those blocks.
Figure~\ref{fig:well_locations} therefore serves as the entry point to the dataset. It shows
where the wells are located, how many labeled samples each well contributes, and how much of
its log suite is missing on average.

\begin{figure}[H]
\centering
\includegraphics[width=0.78\linewidth]{figures/fig01_well_locations_v3.pdf}
\ShortCenteredFigCaption{Spatial distribution of FORCE 2020 wells.}{fig:well_locations}
\end{figure}

In Figure~\ref{fig:well_locations}, marker color represents the per-well mean missing rate,
and marker size represents the number of labeled samples in each well. The figure shows that
data quality is not uniformly distributed across the basin. Nearby wells can share similar
missingness patterns, while different blocks exhibit different acquisition regimes. This
Consequently, a held-out well may represent a different acquisition regime with a new
combination of available and unavailable curves.

The second source of difficulty is the label distribution. Even when curves are present, a
model trained to optimize aggregate performance can be dominated by the most frequent
lithologies. Figure~\ref{fig:lithology_longtail} summarizes this imbalance.

\begin{figure}[H]
\centering
\includegraphics[width=0.82\linewidth]{figures/fig02_lithology_distribution_line.pdf}
\ShortCenteredFigCaption{Long-tailed lithology class distribution.}{fig:lithology_longtail}
\end{figure}

Figure~\ref{fig:lithology_longtail} shows a long-tailed distribution spanning roughly four
orders of magnitude. Shale contains $877{,}043$ samples, accounting for 61.3\% of all valid
labels, whereas Basement contains only 141 samples, or about 0.01\%. The resulting
head-to-tail ratio is approximately $6220{:}1$, and the distribution has a Gini coefficient
of about 0.77. This imbalance explains why ordinary classifiers tend to over-optimize thick,
head-class intervals while under-learning rare lithologies and thin stratigraphic events.

In addition to the spatial heterogeneity in Figure~\ref{fig:well_locations} and the label
imbalance in Figure~\ref{fig:lithology_longtail}, curve availability varies jointly across
wells and log families. Figure~\ref{fig:missingness_matrix} summarizes this structure as a
well-by-curve missing-fraction matrix.

\begin{figure}[H]
\centering
\includegraphics[width=0.86\linewidth]{figures/fig03_missingness_vAB_merged.pdf}
\ShortCenteredFigCaption{Well-by-curve missing-fraction matrix.}{fig:missingness_matrix}
\end{figure}

Figure~\ref{fig:missingness_matrix} groups the curves into the petrophysical families Gamma,
Density/Neutron, Sonic, Resistivity, SP, Caliper, and Drilling/mud. Wells are sorted by their
mean missing rate, while the marginal bars summarize per-curve and per-well missingness. The
block patterns show that missingness is structured by both well and curve family. Entire log
families, such as Sonic or Resistivity, can be absent in specific wells, whereas other curve
families remain available. The observed missingness is therefore inconsistent with a simple
missing-completely-at-random assumption. It reflects tool-string selection, acquisition
period, and operational conditions, which may also vary with the geological targets of
different wells. The representation layer in Section~\ref{sec:repr} encodes missingness and
curve-suite identity directly rather than discarding these patterns during preprocessing.

Finally, robust lithology prediction also requires a geology-aware evaluation target. The
FORCE 2020 benchmark provides an official $12{\times}12$ penalty matrix, shown in
Figure~\ref{fig:penalty}. This matrix assigns small costs to sedimentologically adjacent
confusions, such as Sandstone versus Sandstone/Shale, and larger costs to distant confusions,
such as Sandstone versus Halite.

\begin{figure}[H]
\centering
\includegraphics[width=0.72\linewidth]{figures/fig2_penalty_matrix.pdf}
\ShortCenteredFigCaption{Geological penalty matrix.}{fig:penalty}
\end{figure}

The penalty matrix introduces an asymmetric evaluation criterion. Entry $A_{ij}$ is the cost
of predicting class $j$ when
the true class is $i$; the diagonal cost is zero, within-environment confusions are relatively
cheap, and cross-environment confusions can cost up to 4. Consequently, posterior argmax is
not necessarily the best geological decision. GeoMERIT uses this matrix at the decision layer
by minimizing expected geological penalty rather than treating all errors as equal.

\subsection{Evaluation Metrics: Three Facets of Robustness}
\label{sec:metrics}
We report five complementary metrics. Weighted F1 measures aggregate classification quality,
whereas Macro F1 gives the unweighted mean over the 12 classes. Boundary F1 is computed only
on samples within one depth step of a lithology change and therefore measures stratigraphic
transition quality. Penalty is the mean of $-A_{y\hat y}$ under the FORCE matrix, with values
closer to zero indicating less severe geological errors. Tail mean F1 is the unweighted mean
F1 over the six rarest classes, namely Tuff, Halite, Coal, Dolomite, Anhydrite, and Basement. The three headline metrics probe three distinct facets
of robust prediction and are not interchangeable. Weighted F1 reflects overall prediction
quality and can rise simply by getting more Shale rows right. Boundary F1 reflects stability
in formation-transition zones and rises only by resolving transitions. Penalty reflects the
geological severity of the errors that remain and improves only by avoiding geologically
distant confusions. A method that improves one facet by sacrificing another is not robust;
we therefore evaluate all three jointly.

\subsection{Held-Out-Well Validation as a Robustness Test}
\label{sec:protocol}
We evaluate exclusively under 10-fold GroupKFold cross-validation with well identity as the
grouping key. This protocol is the experimental condition that tests robustness, not part of
the method's definition. GeoMERIT's goal is not to adapt one specific well to another, but to
remain stable under the curve missingness, acquisition heterogeneity, long-tail structure,
and cost asymmetry that any previously unseen well may present, and holding out whole wells
is how that stability is measured. Random row-wise splits place adjacent depth samples from
the same well on both sides of the train/test boundary, allowing models to exploit
well-specific acquisition signatures and local stratigraphic continuity; such splits estimate
interpolation performance within already-logged wells and systematically overstate practical
performance~\cite{rosenblatt2024,roberts2017cv,mila2022nndm}. All feature construction that
uses depth context---window statistics, gradients, boundary masks---is computed strictly
within each well sequence, so no information crosses well boundaries at any stage.

This protocol deliberately exposes a rare-class limitation. For classes concentrated in a few
wells, some training folds contain no positive examples. Basement (141 rows) is the extreme
case. This is not an artifact of the split but a property of the prediction problem itself---a
new well may contain a lithology the training fleet never sampled---so we report Basement
separately and discuss open-set treatment as the appropriate future remedy
(Section~\ref{sec:failures}).

\subsection{Statistical Reporting}
\label{sec:stats}
Headline values are pooled out-of-fold (OOF) metrics over all $1{,}431{,}383$ predictions.
Because pooled values hide fold-to-fold variation, we additionally report the mean
$\pm$ standard deviation across the 10 folds. We also provide an exploratory one-sided
one-sample Wilcoxon signed-rank test~\cite{wilcoxon1945} of the 10 fold values against the
published pooled RFE+XGBoost reference (alternative: GeoMERIT exceeds the reference).
Because fold-matched baseline values are unavailable, this test is descriptive and is not
interpreted as a paired comparison between methods (Section~\ref{sec:rq1}).

\section{GeoMERIT: A Four-Layer Robust Prediction Framework}
\label{sec:method}
The GeoMERIT framework comprises four components: a missingness-aware representation, a
weighted learning objective, structured posterior estimation, and a geological-cost-guided
decision rule (Table~\ref{tab:layers}, Figure~\ref{fig:framework}). Let $x_i$ denote the log
response vector at depth sample $i$, $w(i)$ its training weight,
$p(y\,|\,\phi(x_i))$ the class posterior for the representation $\phi(x_i)$, and $A$ the
FORCE penalty matrix. The representation retains data-quality and curve-suite information
(Section~\ref{sec:repr}). Model fitting uses class-, boundary-, and confidence-aware sample
weights (Section~\ref{sec:learning}), and posterior estimates combine a model ensemble, a
geological hierarchy, and tail experts (Section~\ref{sec:posterior}). After long-tail logit
adjustment, the predicted label is
\begin{equation}
\hat y_i = \arg\min_{j\in\{1,\ldots,12\}} \sum_{k=1}^{12} A_{kj}\;\hat p\!\bigl(y=k\,\big|\,\phi(x_i)\bigr),
\label{eq:bayes}
\end{equation}
as described in Section~\ref{sec:decision}. Equation~(\ref{eq:bayes}) is
the classical Bayes decision rule under an asymmetric loss~\cite{berger1985decision}. It
reduces to argmax when $A$ is the 0--1 loss and departs from argmax exactly where
geological costs are asymmetric.

\begin{table}[H]
\caption{The four robustness layers of GeoMERIT. Each layer counters one condition of real
well-log data that destabilizes conventional lithology classifiers.\label{tab:layers}}
\begin{tabularx}{\textwidth}{lXXl}
\toprule
\textbf{Layer} & \textbf{Real-data condition} & \textbf{GeoMERIT mechanism} & \textbf{Section}\\
\midrule
Representation & Structured curve availability, anomalous responses, inter-well acquisition shift &
Sentinel encoding, missing/outlier indicators, within-well robust normalization, multi-scale
depth context, curve-suite fingerprint & \ref{sec:repr}\\
Learning & Long-tailed classes, ambiguous boundary samples, uncertain labels &
Multiplicative class $\times$ boundary $\times$ confidence sample weighting & \ref{sec:learning}\\
Posterior & Joint 12-class estimation under severe imbalance &
Three-backend GBDT geometric fusion, coarse-to-fine geological hierarchy, one-vs-rest tail experts &
\ref{sec:posterior}\\
Decision & Asymmetric geological misclassification cost ignored by argmax &
Logit adjustment plus Bayes-risk decoding over the FORCE penalty matrix & \ref{sec:decision}\\
\bottomrule
\end{tabularx}
\end{table}

\begin{figure}[H]
\centering
\includegraphics[width=\linewidth]{figures/geo_v2.pdf}
\caption{The GeoMERIT four-layer robust design. Each real-data challenge (left) is countered
by one mechanism layer (right). The representation layer encodes structured curve availability
and anomalous responses. The learning layer increases the influence of tail classes and
boundary samples, while the posterior layer combines a three-backend ensemble, a
coarse-to-fine geological hierarchy, and tail experts. The decision layer replaces argmax
with expected-cost minimization over the FORCE penalty matrix. Evaluation uses strict
10-fold GroupKFold by well.}
\label{fig:framework}
\end{figure}

\subsection{Missingness-Aware Geologic State Representation (Representation Layer)}
\label{sec:repr}
Missingness in well-log collections is often structured and need not satisfy a
missing-completely-at-random assumption. Curve availability depends on factors such as the
tool string, acquisition period, and borehole conditions; consequently, entire log families
may be absent over a well (Figure~\ref{fig:missingness_matrix}). Conventional imputation can
introduce reconstruction error and remove information about the acquisition context. GeoMERIT
therefore retains the missingness pattern explicitly in the input representation.

Concretely, for every curve $c$ we detect missing entries (NaN/Inf) and anomalies (within-well
robust $z$-score $|z|>5$) and produce a sentinel-encoded copy $c_\text{enc}$ in which missing
values become $-999$ and outliers $-777$, alongside binary indicators $c_\text{is\_missing}$
and $c_\text{is\_outlier}$; per-row aggregates of the missing count, missing ratio, outlier
count, and outlier ratio summarize the data quality at each depth. In a GBDT, sentinel values
can be assigned to dedicated branches, allowing split rules to condition on both observed
responses and curve availability without a separate imputation model.

Within each well $w$ and curve $c$ we add a robust normalization
\begin{equation}
z_{w,c}(i) = \frac{c(i)-\operatorname{median}_w(c)}{\operatorname{MAD}_w(c)+\epsilon},\qquad \epsilon=10^{-9},
\label{eq:zscore}
\end{equation}
with a standard-deviation fallback when $\operatorname{MAD}_w(c)=0$. Both the absolute
encoded value and the well-relative value are retained. The absolute response separates
lithologies with distinctive global signatures (e.g., Halite), while the well-relative offset
partly compensates for calibration and borehole-environment drift between wells.

Local depth context is captured by window statistics and stabilized gradients. For the eight
primary curves (GR, RHOB, NPHI, PEF, DTC, RDEP, RMED, SP) and radii $r\in\{2,5,10\}$ we
compute window mean, standard deviation, minimum, maximum, and endpoint difference
($8\!\times\!3\!\times\!5=120$ features), plus the depth-normalized gradient
\begin{equation}
\nabla c(i) = \frac{c(i)-c(i-1)}{\text{DEPTH}(i)-\text{DEPTH}(i-1)+\epsilon}.
\label{eq:grad}
\end{equation}
All context features are computed strictly within each well. Four diagnostic combinations
from the classical interpretation playbook (density--porosity separation $\text{RHOB}-\text{NPHI}$;
the ratios $\text{DTC}/|\text{RHOB}|$, $\text{GR}/\log(1+\text{RDEP})$, and
$\text{PEF}/|\text{RHOB}|$) inject petrophysical priors for separating carbonates from
clastics~\cite{rider2011,doveton2014}.

Finally, the per-well curve-availability bit string is hashed into a categorical
\texttt{curve\_suite\_id}---a single feature that captures vendor-, era-, and tool-driven
variation---and well-head coordinates, distance to the dataset centroid, NPD GROUP and
FORMATION labels, and casing-depth differences complete the state vector $\phi(x_i)$. Missing
stratigraphic labels are encoded as the explicit category ``Unknown'' rather than dropped.

\subsection{Long-Tail- and Boundary-Enhanced Risk Weighting (Learning Layer)}
\label{sec:learning}
The class distribution is strongly imbalanced: Shale accounts for approximately $61\%$ of
the samples, whereas several lithologies have limited support. Under uniform sample
weighting, the empirical loss is therefore dominated by majority-class and
formation-interior samples. To increase the contribution of rare classes and boundary
intervals during fitting, each training sample is assigned the multiplicative weight
\begin{equation}
w(i) = w_\text{class}\!\bigl(y(i)\bigr)\cdot w_\text{bnd}(i)\cdot w_\text{conf}(i).
\label{eq:weight}
\end{equation}
The class term uses the effective-number-of-samples
formulation~\cite{cui2019classbalanced},
\begin{equation}
w_\text{class}(k) = \min\!\left(\frac{1-\rho}{1-\rho^{n_k}},\; w_\text{max}\right),\qquad \rho=0.9995,\quad w_\text{max}=8.0,
\label{eq:effn}
\end{equation}
where $n_k$ is the class support in the training fold; weights are rescaled to unit mean.
Equation~(\ref{eq:effn}) functions as a saturation-aware inverse frequency. Once a class has
enough samples that new ones are redundant, additional support stops reducing its weight,
while the cap $w_\text{max}$ prevents the rarest classes from destabilizing training.
Boundary samples (the label differs from an immediate depth neighbor) receive weight $1.2$,
whereas interior samples receive weight $0.7$. The confidence term maps the FORCE
interpreter-confidence field
$1/2/3$ to $1.0/0.8/0.6$, discounting samples the original labelers themselves marked as
uncertain. We deliberately avoid synthetic oversampling~\cite{chawla2002smote} because interpolated
log rows can violate depth continuity and produce physically impossible curve combinations.

\subsection{Structured Posterior Modeling: Ensemble, Hierarchy, and Tail Experts (Posterior Layer)}
\label{sec:posterior}
A single multiclass model estimates all 12 lithologies jointly despite large differences in
class support and separability. Under this setting, rare classes contribute relatively few
training examples, while adjacent lithologies such as Limestone and Chalk require finer
discrimination than well-separated classes. GeoMERIT retains a multiclass model and augments
its posterior estimates with three complementary components: backbone ensembling, a
geological hierarchy, and one-vs-rest tail experts.

The first mechanism replaces the lone classifier with three complementary backbones trained
per fold---LightGBM~\cite{ke2017lightgbm} on the full training partition ($\approx\!1.28$M rows),
XGBoost~\cite{chen2016xgboost} on a 750{,}000-row subsample, and
CatBoost~\cite{prokhorenkova2018catboost} on a 300{,}000-row subsample, both subsamples
preserving all tail rows. CatBoost uses $\texttt{classes\_count}=12$ to maintain a consistent
12-class output in folds without examples of every class. The backbones differ in their
inductive biases (leaf-wise growth with
native categorical support; depth-wise regularized histogram splitting; ordered boosting with
target statistics). Their posteriors are fused geometrically in log space as
\begin{equation}
\hat p(y=k\,|\,\phi)\propto\exp\!\left(\sum_m \alpha_m\log p^{(m)}(y=k\,|\,\phi)\right),\qquad
\boldsymbol\alpha=(1.0,\,0.8,\,0.8).
\label{eq:geomean}
\end{equation}
The geometric mean assigns a low fused probability to a class when any backend assigns it a
near-zero probability, thereby providing a more conservative pooling rule than arithmetic
averaging.

The second mechanism reorganizes the label space according to geological similarity. The 12
classes are mapped to six super-classes that mirror interpreter taxonomy
(Table~\ref{tab:superclass}); a coarse classifier predicts the super-class, per-group experts
predict the fine class within it, and the hierarchical posterior
\begin{equation}
p_\text{c2f}(y=k\,|\,\phi) = p\!\bigl(\text{super}(k)\,|\,\phi\bigr)\cdot p\!\bigl(k\,|\,\phi,\text{super}(k)\bigr)
\label{eq:c2f}
\end{equation}
is mixed with the multiclass posterior using weight $\mu=0.20$ (validation grid). The
hierarchical component combines group-level and within-group estimates. For example, Chalk
and Limestone are first assigned to the carbonate group and are then distinguished using the
corresponding fine classifier.

\begin{table}[H]
\caption{Geological super-classes used by the coarse-to-fine cascade.\label{tab:superclass}}
\begin{tabularx}{\textwidth}{lX}
\toprule
\textbf{Super-class} & \textbf{Fine classes}\\
\midrule
Clastic & Sandstone, Sandstone/Shale, Shale\\
Carbonate & Limestone, Chalk, Dolomite, Marl\\
Evaporite & Halite, Anhydrite\\
Coal & Coal\\
Basement & Basement\\
Tuff & Tuff\\
\bottomrule
\end{tabularx}
\end{table}

The third mechanism adds dedicated rare-class discriminators that work outside the 12-way
competition. For each tail class $k\in\mathcal T=\{\text{Tuff, Halite,
Coal, Dolomite, Anhydrite, Basement}\}$ a one-vs-rest LightGBM expert with $5\times$
up-weighted positives produces a score $s_k$. When the expert is confident, the corresponding
posterior entry is gated upward:
\begin{equation}
\tilde p(y=k\,|\,\phi) = \hat p(y=k\,|\,\phi)\!\left(1+\gamma_k\,\mathbf 1\bigl[s_k>\theta_k\bigr]\right),
\label{eq:gate}
\end{equation}
followed by renormalization. The gate restores posterior mass that the long tail suppressed,
but only where a dedicated binary discriminator---trained without the 12-way
competition---produces a score above the selected threshold. A single global pair
$(\gamma,\theta)$ serves all tail
classes, with a separate, stricter Dolomite pair $(\gamma_\text{dol},\theta_\text{dol})$
because Dolomite trades off directly against the penalty (Section~\ref{sec:rq5}).

\subsection{Geological-Cost-Guided Decision Layer (Decision Layer)}
\label{sec:decision}
Posterior argmax does not account for asymmetric misclassification costs, and a posterior
estimated from an imbalanced training fold reflects the corresponding class prior. The
decision layer first applies logit
adjustment~\cite{menon2021logit},
\begin{equation}
\hat p_\text{LA}(y=k\,|\,\phi)\propto \tilde p(y=k\,|\,\phi)\cdot \pi_k^{-\tau},
\label{eq:logitadj}
\end{equation}
where $\pi_k$ is the training-fold prior and $\tau\geq 0$ a temperature. This adjustment
reduces the influence of the empirical head-class prior. The adjusted posterior then enters the
Bayes-risk rule of Equation~(\ref{eq:bayes}), so the final label is the one with the lowest
expected geological cost rather than the highest probability. The full decision layer has
five interpretable knobs---$\tau$, the global gate $(\gamma,\theta)$, and the Dolomite gate
$(\gamma_\text{dol},\theta_\text{dol})$---searched on a small validation grid (192
configurations), yielding $\tau=0.2$, $\gamma=1.2$, $\theta=0.35$,
$\gamma_\text{dol}=8.0$, $\theta_\text{dol}=0.05$. The Dolomite setting uses a low activation
threshold together with strong posterior amplification; its effect is therefore assessed
empirically in the sensitivity analysis rather than interpreted as a high-confidence gate.
These parameters can be retuned without retraining the underlying classifiers.

An optional Viterbi pass over each well's depth sequence~\cite{viterbi1967,lafferty2001crf,sarawagi2004semimarkov}
can smooth the decoded labels with an empirical transition matrix; it is examined in
ablation but excluded from headline numbers to keep the decision layer purely pointwise and
tree-only.

\section{Experimental Setup}
\label{sec:setup}

\subsection{Compared Methods and Fairness}
\label{sec:compared}
The primary references are Baseline XGBoost (no feature engineering or weighting),
RFE+XGBoost, and the RFE+\{Random Forest, CatBoost, LightGBM\} variants, all as published
in~\cite{chen2026rfe} and all evaluated on FORCE 2020 under a comparable well-grouped 10-fold
protocol. Their values are used as reported at the operating points selected by the original
authors. GeoMERIT's decision-layer grid was searched on validation folds under the same
GroupKFold protocol used for reporting.

\subsection{Implementation and Scale}
\label{sec:impl}
All experiments ran on a single 112-core CPU machine with 251~GB RAM (no GPU). The software
stack is LightGBM, XGBoost, CatBoost, scikit-learn~\cite{pedregosa2011sklearn}, NumPy,
Pandas, and PyArrow. Each fold trains the three backends, the coarse-to-fine cascade, and
six tail experts; full 10-fold training takes about 3~h and the decoding-parameter search
about 1~h. The feature table occupies $\approx\!12$~GB in Parquet form and peak training
memory is $\approx\!90$~GB.

\section{Results}
\label{sec:results}
The evaluation contains five analyses. Section~\ref{sec:rq1} reports overall, boundary, and
cost-sensitive performance under held-out-well validation together with fold-level
variability. Section~\ref{sec:rq2} compares the framework with published feature-engineering
and tree-backbone references and examines a same-stack decoder ablation.
Section~\ref{sec:rq3} reports performance after stratifying held-out wells by missingness,
and Section~\ref{sec:rq4} presents class-level and well-track diagnostics.
Section~\ref{sec:rq5} evaluates sensitivity to the decision-layer hyperparameters.

\subsection{Main Comparison Under Held-Out-Well Validation}
\label{sec:rq1}
The primary comparison evaluates three complementary outcomes under 10-fold GroupKFold by
well: Weighted F1 for aggregate classification quality, Boundary F1 for formation-transition
samples, and Penalty for the geological severity of errors. Relative to the published
RFE+XGBoost reference, GeoMERIT raises Weighted F1 from $0.7272$ to $0.7526$
($+0.0254$; $+3.5\%$), Boundary F1 from $0.4100$ to $0.4567$ ($+0.0467$; $+11.4\%$), and
Penalty from $-0.6289$ to $-0.5886$ ($+0.0403$; $+6.4\%$) (Table~\ref{tab:main},
Figure~\ref{fig:main}). The simultaneous improvement indicates that the gain in aggregate
performance is not accompanied by lower boundary performance or a higher geological error
cost. The tail mean F1 is $0.5061$, compared with approximately $0.31$ for the single-model
baseline level.

\begin{table}[H]
\caption{Main comparison under 10-fold GroupKFold by well. Reference values from~\cite{chen2026rfe};
GeoMERIT values are pooled out-of-fold over all $1{,}431{,}383$ samples. The third row
(``base ensemble $+$ Bayes'') uses the representation, learning, and posterior layers of
GeoMERIT but decodes with plain Bayes risk (no logit adjustment, no tail gates); contrasting
it with the full row isolates logit adjustment and tail gating within the decision layer. It
does not isolate Bayes-risk decoding from argmax. $\Delta$ rows give the change relative to
RFE+XGBoost.\label{tab:main}}
\begin{tabularx}{\textwidth}{Xccccc}
\toprule
\textbf{Method} & \textbf{Weighted F1} & \textbf{Macro F1} & \textbf{Boundary F1} & \textbf{Penalty} & \textbf{Tail F1}\\
\midrule
Baseline XGBoost~\cite{chen2026rfe} & 0.6640 & -- & 0.3510 & $-0.7800$ & --\\
RFE + XGBoost~\cite{chen2026rfe} & 0.7272 & -- & 0.4100 & $-0.6289$ & --\\
GeoMERIT, base ensemble + Bayes decoding & 0.7446 & 0.5038 & 0.4423 & $-0.5885$ & 0.4737\\
\textbf{GeoMERIT, full decision layer} & \textbf{0.7526} & \textbf{0.5315} & \textbf{0.4567} & $\mathbf{-0.5886}$ & \textbf{0.5061}\\
\midrule
$\Delta$ absolute vs.\ RFE+XGBoost & $+0.0254$ & -- & $+0.0467$ & $+0.0403$ & --\\
$\Delta$ relative vs.\ RFE+XGBoost & $+3.5\%$ & -- & $+11.4\%$ & $+6.4\%$ & --\\
\bottomrule
\end{tabularx}
\end{table}

\begin{figure}[H]
\centering
\includegraphics[width=\linewidth]{figures/fig4_main_results.pdf}
\caption{GeoMERIT improves aggregate, boundary-sensitive, and cost-sensitive performance
simultaneously under held-out-well validation. Panels show Weighted F1, Boundary F1, and
Penalty for Baseline XGBoost, RFE+XGBoost~\cite{chen2026rfe}, and GeoMERIT; annotations give
the gain over the RFE reference. The Penalty axis is signed so that closer to zero is
better.}
\label{fig:main}
\end{figure}

Pooled metrics do not characterize variation among held-out well groups. Table~\ref{tab:perfold}
therefore reports fold-level dispersion and exploratory one-sample tests against the pooled
RFE+XGBoost reference. Boundary F1 exceeds the reference in 9 of 10 folds ($p=0.002$),
Weighted F1 in 7 of 10 folds ($p=0.080$), and Penalty in 6 of 10 folds ($p=0.138$).
However, the published reference is a pooled point estimate rather than a matched set of
per-fold results. These tests consequently do not provide a paired comparison between methods
and are not used to claim inferential superiority. Figure~\ref{fig:perfold} reports the
underlying fold values to show the magnitude of geological heterogeneity across splits.

\begin{table}[H]
\caption{Fold-level dispersion of GeoMERIT (mean $\pm$ SD over 10 GroupKFold folds), pooled
OOF value, and exploratory one-sided one-sample Wilcoxon signed-rank test against the pooled
RFE+XGBoost reference value. Because fold-matched reference values are unavailable, the test
is descriptive rather than a paired comparison. Pooled Macro F1 exceeds the fold mean because
several tail classes are concentrated in a few wells. When a held-out fold contains no samples
of a class, the fold-level calculation assigns that class an F1 of zero, which lowers the mean
of fold-level Macro F1 values but not the pooled value.\label{tab:perfold}}
\begin{tabularx}{\textwidth}{Xcccc}
\toprule
\textbf{Metric} & \textbf{Fold mean $\pm$ SD} & \textbf{Pooled OOF} & \textbf{Folds $>$ ref.} & \textbf{Wilcoxon $p$}\\
\midrule
Weighted F1 & $0.7521\pm 0.0477$ & $0.7526$ & 7/10 & $0.080$\\
Macro F1 & $0.4410\pm 0.0842$ & $0.5315$ & --- & ---\\
Boundary F1 & $0.4567\pm 0.0250$ & $0.4567$ & 9/10 & $\mathbf{0.002}$\\
Penalty & $-0.5884\pm 0.0937$ & $-0.5886$ & 6/10 & $0.138$\\
\bottomrule
\end{tabularx}
\end{table}

\begin{figure}[H]
\centering
\includegraphics[width=\linewidth]{figures/fig5_perfold_forest.pdf}
\caption{Per-fold metrics across the 10 GroupKFold folds. Each row is one held-out fold;
filled markers exceed the pooled RFE+XGBoost reference (dashed line) and open markers fall
below. The diamond row gives mean $\pm$ SD. Fold-to-fold spread reflects the geological
heterogeneity of the held-out well groups, motivating the reporting of dispersion in addition
to pooled point estimates.}
\label{fig:perfold}
\end{figure}

\subsection{Analysis of Performance Gains}
\label{sec:rq2}
To assess whether the improvement can be explained solely by the choice of base classifier,
we compare GeoMERIT with the published RFE ablation sequence, examine results across tree
backbones, and evaluate a same-stack decoder contrast. In the published sequence
of~\cite{chen2026rfe} (Table~\ref{tab:ladder}),
quality-aware weighting and feature engineering each close part of the gap from the $0.6635$
baseline to $0.7272$, and GeoMERIT's added layers---the three-backend geometric ensemble, the
geological hierarchy, the tail experts, and Bayes-risk decoding---close a further $+0.0254$
Weighted F1 and $+0.0403$ Penalty on top of them. Because those rungs were not re-run inside
our implementation, the final increment is a cumulative cross-study comparison rather than a
per-component ablation. Sections~\ref{sec:rq3}--\ref{sec:rq4} provide targeted diagnostics
for the representation and learning objectives, while the same-stack comparison below
isolates logit adjustment and tail gating.

\begin{table}[H]
\caption{Comparison with published RFE configurations from~\cite{chen2026rfe}. The
GeoMERIT increment is a cumulative cross-study comparison and should not be interpreted as a
per-layer causal decomposition.\label{tab:ladder}}
\begin{tabularx}{\textwidth}{Xcc}
\toprule
\textbf{Configuration} & \textbf{Weighted F1} & \textbf{Penalty}\\
\midrule
Baseline XGBoost (no FE, no weighting) & 0.6635 & $-0.7804$\\
Baseline $+$ weighting & 0.7230 & $-0.6346$\\
Baseline $+$ feature engineering & 0.7259 & $-0.6327$\\
RFE $+$ XGBoost (FE $+$ weighting) & 0.7272 & $-0.6289$\\
\textbf{GeoMERIT} ($+$ ensemble, hierarchy, tail experts, Bayes-risk decoding) & \textbf{0.7526} & $\mathbf{-0.5886}$\\
\bottomrule
\end{tabularx}
\end{table}

The backbone comparison provides a second reference (Table~\ref{tab:backbones}). RFE-style
features improve each reported single-model variant, consistent with~\cite{chen2026rfe}, and
each remains below GeoMERIT on all three metrics. The difference between the best reported
single model (RFE+LightGBM at $0.7092$) and GeoMERIT ($0.7526$) is consistent with the
additional use of geometric fusion, the hierarchy, tail gating, and cost-aware decoding,
although the cross-study design does not identify their individual contributions.

\begin{table}[H]
\caption{Tree-backbone reference rungs (from~\cite{chen2026rfe}) versus the GeoMERIT
ensemble. Every reported single-backbone RFE variant remains below the GeoMERIT operating
point on all three metrics.\label{tab:backbones}}
\begin{tabularx}{\textwidth}{Xccc}
\toprule
\textbf{Model} & \textbf{Weighted F1} & \textbf{Boundary F1} & \textbf{Penalty}\\
\midrule
Random Forest & 0.6918 & 0.3840 & $-0.6776$\\
RFE $+$ Random Forest & 0.7051 & 0.3920 & $-0.6467$\\
CatBoost & 0.6501 & 0.3510 & $-0.7895$\\
RFE $+$ CatBoost & 0.6866 & 0.3725 & $-0.7156$\\
LightGBM & 0.6550 & 0.3768 & $-0.8080$\\
RFE $+$ LightGBM & 0.7092 & 0.4147 & $-0.6691$\\
\textbf{GeoMERIT ensemble} & \textbf{0.7526} & \textbf{0.4567} & $\mathbf{-0.5886}$\\
\bottomrule
\end{tabularx}
\end{table}

Holding the representation, learning, posterior estimates, and Bayes-risk rule fixed, the
third and fourth rows of Table~\ref{tab:main} isolate the addition of logit adjustment and
tail gating. These components add $+0.0080$ Weighted F1, $+0.0277$
Macro F1, $+0.0144$ Boundary F1, and $+0.0324$ Tail F1 at essentially unchanged Penalty
($-0.5885\to-0.5886$). This same-stack contrast directly measures the combined effect of the
two posterior-adjustment mechanisms. It does not, however, isolate Bayes-risk decoding because
both rows use that rule. The remaining difference relative to the published RFE result reflects
the representation, learning, posterior, and Bayes-risk components jointly and cannot be
decomposed from the reported cross-study values.

\subsection{Robustness on Feature-Deficient Wells}
\label{sec:rq3}
To evaluate performance under incomplete curve suites, we stratify held-out wells by mean
per-row missingness and compare feature-complete and feature-deficient subsets at several
thresholds (Figure~\ref{fig:robust}). Each subset is annotated with its number of wells so
that results from sparsely populated subsets can be identified. The $38$-well deficient
bucket ($0.1/0.4$ split) attains Weighted F1
$0.7103$ and Boundary F1 $0.4408$---above the full-fleet RFE+XGBoost boundary reference of
$0.4100$ despite operating only on incomplete wells. The $8$-well deficient bucket
($0.2/0.5$) attains Weighted F1 $0.7555$, exceeding the full-fleet reference of $0.7272$
outright. At the $0.3/0.6$ threshold, the deficient subset contains only one well with
$11{,}315$ rows and has substantially lower performance. The populated subsets indicate that
the model can retain useful performance under incomplete curve suites, while the single-well
result also shows that robustness is not uniform across acquisition regimes. Because the
subsets differ in geological composition, this stratification is a descriptive stress test
rather than a controlled ablation of individual representation features.

\begin{figure}[H]
\centering
\includegraphics[width=\linewidth]{figures/fig6_missingness_robustness.pdf}
\caption{Missingness-stratified robustness under held-out-well validation. At each threshold,
wells are split into feature-complete (solid blue) and feature-deficient (dashed orange)
subsets by mean per-row missing rate; annotations give bucket sizes in wells. The gray
dashed line is the pooled full-fleet RFE+XGBoost reference. Deficient subsets containing 38
and 8 wells remain near or above the reference, whereas the $0.3/0.6$ deficient subset
contains only one well and should be interpreted accordingly.}
\label{fig:robust}
\end{figure}

\subsection{Per-Class Diagnosis, Tail Behavior, and Well Tracks}
\label{sec:rq4}
Because Weighted F1 is strongly influenced by Shale, aggregate performance may obscure weak
tail classes and boundary behavior. We therefore report per-class out-of-fold precision,
recall, and F1 (Table~\ref{tab:perclass}). Figure~\ref{fig:classdiag} relates class-level F1
to support (panel~a), summarizes the destinations of misclassified samples (panel~b), and
reports performance for the six classes targeted by the tail experts (panel~c).

\begin{table}[H]
\caption{Per-class out-of-fold metrics for the 12 FORCE 2020 lithology classes (pooled over
10 folds). Support is the number of labeled rows. Figure~\ref{fig:classdiag}a visualizes the
F1 column jointly with support; precision and recall are reported here
only.\label{tab:perclass}}
\begin{tabularx}{\textwidth}{Xrccc}
\toprule
\textbf{Class} & \textbf{Support} & \textbf{Precision} & \textbf{Recall} & \textbf{F1}\\
\midrule
Shale & 877{,}043 & 0.833 & 0.912 & 0.870\\
Sandstone & 207{,}704 & 0.768 & 0.723 & 0.745\\
Sandstone/Shale & 180{,}820 & 0.442 & 0.285 & 0.347\\
Limestone & 69{,}498 & 0.636 & 0.659 & 0.647\\
Marl & 41{,}038 & 0.382 & 0.368 & 0.375\\
Tuff & 17{,}431 & 0.663 & 0.742 & 0.700\\
Halite & 14{,}712 & 0.968 & 0.983 & 0.975\\
Chalk & 14{,}043 & 0.815 & 0.228 & 0.357\\
Coal & 4{,}754 & 0.649 & 0.692 & 0.670\\
Dolomite & 2{,}391 & 0.040 & 0.041 & 0.040\\
Anhydrite & 1{,}808 & 0.606 & 0.704 & 0.651\\
Basement & 141 & 0.000 & 0.000 & 0.000\\
\bottomrule
\end{tabularx}
\end{table}

\begin{figure}[H]
\centering
\includegraphics[width=\linewidth]{figures/fig7_class_diagnosis.pdf}
\caption{Class-level diagnostic results for GeoMERIT.
(\textbf{a})~Per-class F1 with marker size proportional to log support (support in
parentheses); classes ordered by geological group with the shared lithology palette.
(\textbf{b})~Row-normalized confusion matrix in the same group order as
Figure~\ref{fig:penalty}; framed blocks show that residual errors concentrate inside the
clastic and carbonate--marl clusters (values below $0.02$ suppressed).
(\textbf{c})~The six tail classes against the $0.65$ recovery goal: four are recovered
(green); Dolomite and Basement remain failure modes (red).}
\label{fig:classdiag}
\end{figure}

Two patterns are evident. First, Halite (F1 $0.975$), Tuff ($0.700$), Coal ($0.670$), and
Anhydrite ($0.651$) exceed $0.65$ despite jointly accounting for less than $3\%$ of the
labels. Their distinctive log responses are consistent with successful rare-class recovery.
Second, the largest residual confusions occur among geologically related classes.
Sandstone/Shale is predicted as Shale for $0.538$ of its samples, and Chalk is predicted as
Limestone for $0.654$. Marl is predicted as Shale and Limestone for $0.408$ and $0.178$,
respectively. Dolomite is predicted as Shale and Limestone for $0.437$ and $0.311$.
Most large off-diagonal entries therefore fall within the low-cost clusters identified by the
penalty matrix (Figure~\ref{fig:penalty}), which is consistent with the cost-aware decision
rule. The tail mean F1 of $0.5061$ is approximately $1.6\times$ the
$\approx\!0.31$ level of the reported single-model baselines.

To complement the pooled analysis, we inspect two held-out wells representing the two
conditions emphasized in Section~\ref{sec:repr}: feature-complete well 16/7-6 and
missingness-heavy 35/9-7, both located within the North Sea study area shown in
Figure~\ref{fig:well_locations}.
Figure~\ref{fig:tracks} shows, for each well, five measured curves (GR, RHOB, NPHI, DTC,
RDEP), the per-row suite missing fraction, and the true/predicted/hit lithology strips on a
common depth axis.

\begin{figure}[H]
\centering
\includegraphics[width=\linewidth]{figures/fig8_well_tracks.pdf}
\caption{Held-out single-well prediction tracks.
(\textbf{a})~Feature-complete well 16/7-6 (point accuracy $75.5\%$): five log curves,
per-row suite missing fraction, and true/predicted/hit strips in the shared lithology
palette; dotted lines bracket the 2200--2400~m carbonate interval, where the chalk section
is decoded as Limestone---the depth-localized expression of the
$\text{Chalk}\to\text{Limestone}$ confusion ($0.654$) of Figure~\ref{fig:classdiag}b.
(\textbf{b})~Missingness-heavy well 35/9-7 (point accuracy $87.6\%$): six of twenty suite
curves (RSHA, SGR, SP, DCAL, MUDWEIGHT, RXO) are entirely absent (gray missing-fraction
track), yet the predicted column reproduces the Shale-dominated sequence and its sandstone
interbeds with only thin boundary errors. Hatched bands mark within-curve gaps.}
\label{fig:tracks}
\end{figure}

The feature-complete well illustrates the class-level confusion pattern. Prediction quality is high along
the clastic section, while the bracketed 2200--2400~m interval---about 1{,}200 rows of Chalk
with thin Marl---is decoded almost entirely as Limestone (interval accuracy $6.0\%$) and
alone accounts for $48\%$ of the well's errors. The error is therefore not diffuse model
noise; it is a depth-contiguous carbonate--marl ambiguity consistent with
Figure~\ref{fig:classdiag}b. The RHOB/NPHI response in this interval lies between the chalk
and limestone training distributions. In the missingness-heavy well, six curves are absent
over the entire interval, point accuracy is $87.6\%$, and most errors occur near thin
transitions. This case illustrates the behavior of the missingness-aware representation at
the well-track scale but does not by itself establish the contribution of an individual feature.

\subsection{Decision-Layer Sensitivity}
\label{sec:rq5}
We evaluate the sensitivity of the decision layer by varying two groups of hyperparameters
while holding the remaining pipeline fixed: the logit-adjustment temperature $\tau$ and the
Dolomite gate $(\gamma_\text{dol},\theta_\text{dol})$ from
Section~\ref{sec:decision} (Figure~\ref{fig:decision}).

\begin{figure}[H]
\centering
\includegraphics[width=\linewidth]{figures/fig9_decision_layer.pdf}
\caption{Decision-layer trade-off under logit adjustment and Dolomite gating.
(\textbf{a})~Sweeping $\tau$ (gates off): the four F1 facets rise as the head-class prior is
flattened while Penalty peaks near $\tau\approx 0.1$ and degrades slowly beyond; the
adopted $\tau=0.2$ (dashed) sits inside the shaded flat region, trading $\le 0.002$ Penalty
for $+0.009$ Macro F1 relative to $\tau=0.1$.
(\textbf{b})~Dolomite-gate grid: each point is one
$(\gamma_\text{dol},\theta_\text{dol})$ setting, colored by $\gamma_\text{dol}$; the red
curve is the Pareto frontier of Dolomite F1 versus Penalty and the star is the adopted
setting $(8.0, 0.05)$. Settings near the maximum Dolomite F1 ($\approx 0.076$) worsen the
Penalty metric by as much as $0.04$.}
\label{fig:decision}
\end{figure}

The $\tau$ sweep is monotone and relatively flat near the adopted value. Between
$\tau=0.1$ and $\tau=0.3$, Weighted F1 varies by $0.0045$ and Penalty by $0.0077$, indicating
limited sensitivity within this interval. The adopted Dolomite-gate setting lies on the
estimated Pareto frontier. Configurations with higher Dolomite F1 have worse (more negative)
Penalty scores,
consistent with overlap between Dolomite and neighboring carbonate classes
(Section~\ref{sec:failures}). Because these adjustments are applied after model fitting, the
trade-off can be examined without retraining the underlying classifiers.

\section{Discussion}
\label{sec:discussion}

\subsection{Interpretation of the Joint Performance Gains}
\label{sec:joint}
Under held-out-well validation, GeoMERIT exceeds the published RFE+XGBoost reference on the
three pooled headline metrics: Weighted F1, Boundary F1, and Penalty. This joint result shows
that the aggregate gain is accompanied by improved boundary and cost-sensitive performance.
The missingness-stratified, class-level, and sensitivity analyses are consistent with the
intended roles of the representation, weighting, posterior, and decision components.
However, the reported experiments do not provide a complete leave-one-component-out ablation.
In particular, the same-stack comparison in Section~\ref{sec:rq2} isolates logit adjustment
and tail gating but not Bayes-risk decoding relative to argmax. The Penalty improvement should
therefore be interpreted as a property of the full framework rather than attributed solely to
the Bayes-risk rule.

\subsection{Geological Structure of Residual Errors}
\label{sec:failures}
The residual errors are concentrated in a small number of classes and can be related to their
geological and sampling characteristics.

\textbf{Dolomite (F1 $=0.040$)} is the largest visible failure, and its cause is mineralogical
rather than statistical. Clean dolomite (PEF $\approx 3.1$, matrix density
$\approx 2.87$~g/cm$^3$) sits squarely inside the carbonate subspace already occupied by
limestone (PEF $\approx 5$, $\approx 2.71$~g/cm$^3$), chalk, and marl, and the gate sweep of
Section~\ref{sec:rq5} showed that overlap to be strong enough that forcing dolomite recall
actively worsens the penalty. This pattern motivates more targeted alternatives, such as a
carbonate-subspace expert, retrieval priors restricted to the same GROUP and FORMATION, or a
dedicated Dolomite-versus-Limestone discriminator trained where the two classes coexist.

\textbf{Basement (F1 $=0.000$)} fails completely, but for a reason that lies in the data rather
than the model. Its 141 rows are concentrated in a few wells where the key curves are
themselves heavily missing, so under GroupKFold several training folds contain no Basement
supervision. This setting may be better addressed by an open-set or anomaly-detection
formulation than by a closed-set multiclass output.

\textbf{Sandstone/Shale ($0.347$) and Chalk ($0.357$)} are weak despite ample support---
Sandstone/Shale alone has $180{,}820$ rows---so here the ceiling is set by the labels, not by
rarity. Sandstone/Shale is by definition a transitional rock between its two parents, and
Chalk's combination of high precision ($0.815$) with low recall ($0.228$) describes a model
that commits to chalk only when the evidence is overwhelming and otherwise defers to limestone.
These results suggest a resolution and label-ambiguity limitation at log scale. Improved
labels or sequence-level decoding over bed intervals may help distinguish these classes.

\subsection{Scope and Limitations}
\label{sec:limits}
Four limitations bound the claims. First, the cross-stack comparison in
Table~\ref{tab:ladder} splices published rungs with our extension, so the four-layer
increment is attributed cumulatively rather than per-layer in that comparison; the
same-stack contrast of Section~\ref{sec:rq2} isolates only logit adjustment and tail gating.
A full in-stack leave-one-component-out study is needed to quantify the remaining
contributions. Second, the effect of Bayes-risk decoding is not isolated from argmax using
identical posterior estimates. Third, the fold-level tests use a pooled published reference
rather than matched baseline folds and are therefore exploratory rather than inferential
comparisons (Table~\ref{tab:perfold}). Fourth, all conclusions are for one basin and one
penalty matrix; the decision layer would inherit any context-dependence of $A$ (e.g., depth-
or formation-conditional costs). Equation~(\ref{eq:bayes}) can accept a different cost matrix
without retraining, but its performance would require validation in the new context.

\subsection{Rationale for a Tree-Based Model Stack}
\label{sec:trees}
GeoMERIT instantiates, in a geological setting, the broader observation that tree ensembles
remain highly competitive on tabular
data~\cite{grinsztajn2022trees,shwartzziv2022tabular,mcelfresh2024benchmark}. GBDTs natively
route the sentinel-encoded missingness states, remain insensitive to the
four-orders-of-magnitude scale differences across curves, use axis-aligned decision surfaces
that mirror the cross-plot logic of petrophysical interpretation, and accept categorical
features such as \texttt{curve\_suite\_id} and FORMATION as first-class inputs. These
properties do not preclude deep-learning extensions. Well-log foundation
models~\cite{qi2025wlfm} could supply
upstream representations that the same weighting, hierarchy, and decoding layers consume,
and agentic workflows~\cite{wang2026geodecider} could invoke GeoMERIT as a calibrated
tool. For the missingness-heavy, long-tailed, and strictly well-grouped setting studied here,
tree ensembles provide a suitable model class for the structured tabular representation.

\section{Conclusions}
\label{sec:conclusion}
We presented GeoMERIT, a missingness-aware and penalty-guided lithology prediction method for
incomplete, heterogeneous, and long-tailed well-log data. The representation encodes
missingness, anomalies, and curve-suite identity. Model fitting applies class, boundary, and
label-confidence weights, while posterior estimation combines a tree ensemble with a
geological hierarchy and tail experts. Final labels minimize expected geological cost after
logit adjustment. On the full FORCE 2020 valid-label set under strict
10-fold GroupKFold validation by well---the protocol that simulates prediction on unseen
wells---GeoMERIT attains Weighted F1 $=0.7526$, Macro F1 $=0.5315$, Boundary F1 $=0.4567$,
Penalty $=-0.5886$, and tail mean F1 $=0.5061$, improving the published RFE+XGBoost
reference by $+3.5\%$ (Weighted F1), $+11.4\%$ (Boundary F1), and $+6.4\%$ (Penalty).
Within the same posterior stack, adding logit adjustment and tail gating increases Weighted
F1 by $0.0080$, Macro F1 by $0.0277$, Boundary F1 by $0.0144$, and Tail F1 by $0.0324$ at
essentially unchanged Penalty. These results support the joint treatment of missingness,
long-tail structure, boundary uncertainty, and geological cost, while the incomplete
component ablation limits causal attribution to individual layers. Future work should evaluate
carbonate-subspace experts for Dolomite, open-set formulations for Basement-like classes,
semi-Markov segment decoding at bed scale, and foundation-model representations under the
same held-out-well protocol.

\vspace{6pt}

\authorcontributions{Conceptualization, methodology, software, validation, formal analysis,
investigation, data curation, writing---original draft preparation, writing---review and
editing, and visualization: A.A. The author has read and agreed to the published version of
the manuscript.}

\funding{This research received no external funding.}

\institutionalreview{Not applicable.}

\informedconsent{Not applicable.}

\dataavailability{The FORCE 2020 dataset is publicly available at
\url{https://doi.org/10.5281/zenodo.4351156}~\cite{bormann2020force}. The code,
configuration files, fold definitions, out-of-fold predictions, decoding grids, and
plotting scripts that reproduce every number, table, and figure in this paper are described
in Appendix~\ref{app:repro}.}

\acknowledgments{The authors thank the FORCE 2020 organizers for releasing the dataset and
penalty matrix that made this study possible.}

\conflictsofinterest{The author declares no conflicts of interest.}

\abbreviations{Abbreviations}{
The following abbreviations are used in this manuscript:\\

\noindent
\begin{tabular}{@{}ll}
GBDT & Gradient-boosted decision tree\\
RFE  & Robust feature engineering\\
OOF  & Out-of-fold\\
MAD  & Median absolute deviation\\
OvR  & One-vs-rest\\
GR   & Gamma ray\\
RHOB & Bulk density\\
NPHI & Neutron porosity\\
PEF  & Photoelectric factor\\
DTC  & Compressional sonic transit time\\
RDEP & Deep resistivity\\
NPD  & Norwegian Petroleum Directorate\\
\end{tabular}}

\appendixtitles{yes}
\appendixstart
\appendix
\section[\appendixname~\thesection]{Reproducibility Artifacts}
\label{app:repro}
All numbers in this paper are regenerated from saved artifacts. Core library:
\texttt{geomerit/features.py} (representation layer), \texttt{weights.py} (learning layer),
\texttt{models.py} (posterior layer: three backends, geometric fusion, coarse-to-fine, tail
experts), \texttt{decode.py} (decision layer), \texttt{metrics.py} (five metrics),
\texttt{cv.py} (GroupKFold protocol). Result files: \texttt{decode\_report.json} (pooled
metrics and per-class report), \texttt{decode\_report.csv} (per-row OOF predictions),
\texttt{decode\_grid\_results.json} (decision-layer grid), \texttt{perfold\_geomerit.csv}
(per-fold metrics; \texttt{scripts/06\_perfold\_stats.py} reproduces the Wilcoxon tests).
Figures: every panel is generated by \texttt{revision\_v3/scripts/make\_figures.py} from
same-named CSV source data; the two case-well curve extracts derive from the public FORCE
2020 release. Best decoding configuration: stage \texttt{global\_tail}, $\tau=0.2$,
$(\gamma,\theta)_\text{global}=(1.2,0.35)$,
$(\gamma,\theta)_\text{dol}=(8.0,0.05)$.

\begin{adjustwidth}{-\extralength}{0cm}
\reftitle{References}
\bibliography{references}
\end{adjustwidth}

\PublishersNote{}
\end{document}
